\section{Software Architecture}
\label{sec:sw}

This section describes the firmware on the \texttt{ltp} branch (commit
\texttt{ca2e48e}), the most recent version. It merges the \texttt{master}
branch with the tack planner and the sensor and stack updates from
\texttt{sensor-upgrade-2}. The firmware is written in C on the STM32 hardware
abstraction layer (HAL),
about 8000 lines of application code across the two cores and a shared
library. It is built with CMake. A single shared configuration header enables
or disables each module at compile time, so any subsystem can be built and
tested alone. The same header selects built-in test modes.

\subsection{Dual-Core Partitioning}
The STM32H755 has a Cortex-M7 and a Cortex-M4 that share memory but run
separate firmware images~\cite{rm0399}. ALAN divides the work by timing
requirement. The CM7 runs the FreeRTOS real-time operating system
(RTOS)~\cite{barry2016} and owns every sensor, both planners and control
loops, both actuators, and the mode logic. The CM4 runs without an operating
system and owns only the radio. Its main loop sleeps on a wait-for-interrupt
instruction and wakes for two events: the radio's DIO0 line, which signals a
finished transmission or a received packet, and timer TIM6, which fires once
per second to start a telemetry transmission. Because radio transfers and
timeouts happen on the other core, they cannot delay a control loop.

\subsection{Inter-Core Communication}
\label{sec:ipc}
The cores exchange data only through a \SI{64}{\kilo\byte} block of SRAM4 at
address \texttt{0x38000000}, which both linker scripts place at the same
address (Fig.~\ref{fig:ipc}). At boot the CM7 configures its memory protection
unit (MPU) to make the region non-cacheable, clears it, and writes a layout
version followed by a magic word (``ALAN'' in ASCII). The CM4 waits for the
magic word and refuses to attach if the version, which combines a revision
number with the structure size, does not match its own. This catches the
common mistake of flashing the cores with different layouts. The block holds
two channels:
\begin{itemize}
\item \emph{Telemetry} (CM7 to CM4): a single latest-value copy of the
telemetry record, guarded by hardware semaphore (HSEM) 1. Only the newest
record matters, so a core that finds the semaphore taken skips that cycle
instead of waiting.
\item \emph{Commands} (CM4 to CM7): a lock-free single-producer,
single-consumer ring of 64 slots of 18 bytes (type, sequence number, and a
16-byte payload). The CM4 writes only the head index and the CM7 only the
tail, and memory barriers order each slot write before its index update.
After each push the CM4 takes and releases HSEM 2, which raises an interrupt
on the CM7 that wakes the command task.
\end{itemize}

\begin{figure}[!t]
\centering
\begin{tikzpicture}[font=\scriptsize, >={Stealth[length=2mm]}]
\node[draw, thick, fill=blue!6, minimum width=16mm, minimum height=28mm, align=center] (m7) at (0,0) {\textbf{CM7}\\[4pt]telemetry\\task\\[8pt]commander\\task};
\node[draw, thick, fill=blue!6, minimum width=16mm, minimum height=28mm, align=center] (m4) at (6.3,0) {\textbf{CM4}\\[4pt]TIM6\\1 Hz tick\\[8pt]LoRa RX\\handler};
\node[draw, fill=gray!10, minimum width=34mm, minimum height=30mm] (s) at (3.15,0) {};
\node[anchor=south, align=center] at (s.north) {SRAM4 at 0x38000000 (non-cacheable)};
\node[draw, fill=white, minimum width=29mm, align=center] (tr) at (3.15,0.7) {telemetry record (73 B)\\guarded by HSEM 1};
\foreach \i in {0,...,7} \draw[fill=white] (1.85+\i*0.325,-0.95) rectangle ++(0.325,0.42);
\node at (3.15,-0.27) {command ring, 64 $\times$ 18 B};
\node[font=\tiny] at (2.0,-1.2) {tail};
\node[font=\tiny] at (4.3,-1.2) {head};
\draw[->] (m7.east |- tr) -- node[above,font=\tiny]{write} (tr.west);
\draw[->] (tr.east) -- node[above,font=\tiny]{read} (m4.west |- tr);
\draw[->] (m4.west |- 0,-0.74) -- node[above,font=\tiny]{push} (4.45,-0.74);
\draw[->] (1.85,-0.74) -- node[above,font=\tiny]{pop} (m7.east |- 0,-0.74);
\draw[->, dashed] (m4.south) -- ++(0,-0.35) -| node[pos=0.25,below]{HSEM 2 doorbell interrupt} (m7.south);
\end{tikzpicture}
\caption{Inter-core communication. Telemetry is a single latest-value record
under a hardware semaphore; commands use a lock-free ring, and a second
semaphore acts as a doorbell that wakes the CM7.}
\label{fig:ipc}
\end{figure}

\subsection{Task Structure and Data Flow}
Table~\ref{tab:tasks} lists the CM7 tasks. Every periodic task waits on
\texttt{vTaskDelayUntil}, so its period is fixed and each PID uses a constant
sample time. Each sensor task publishes its latest result to a single-slot
FreeRTOS queue that is overwritten on every update. Every result carries a
timestamp, a validity code (for example valid, calibrating, window filling,
or untrusted), and a sequence number, so consumers can reject stale or bad
data instead of acting on it. Priorities follow consequence. The commander
runs highest so a mode change takes effect at once. The arbiters and the
telemetry task run above the sensor drivers, whose I\textsuperscript{2}C and
UART reads busy-wait, so servo outputs and the SRAM4 publish are never
starved. The planners run lowest.

\begin{table}[!t]
\caption{CM7 FreeRTOS Tasks}
\label{tab:tasks}
\centering
\begin{tabular}{@{}llr@{}}
\toprule
Task & Priority & Rate (Hz) \\
\midrule
Commander (modes, commands) & High         & event, 10 \\
Telemetry packing           & Above normal & 1 \\
Rudder arbiter              & Above normal & 20 \\
Tail arbiter                & Above normal & 20 \\
Heading PID                 & Normal       & 10 \\
Sail PID                    & Normal       & 10 \\
IMU (BNO055)                & Normal       & 50 \\
Sail encoder (AS5600)       & Normal       & 10 \\
Orientation fusion          & Normal       & on IMU data \\
Wind vane                   & Normal       & 5 \\
GPS                         & Low          & 1 \\
Battery                     & Low          & 2 \\
Live sail planner (LSP)     & Low          & 2 \\
Live tack planner (LTP)     & Normal       & 2 \\
Serial console (debug)      & Low          & on demand \\
\bottomrule
\end{tabular}
\end{table}

Angles need care when averaged. The encoder jumps from 4095 to 0 counts once
per turn, and the wind vane does the same at \SI{180}{\degree}, so a plain
average of readings that straddle the seam points in the opposite direction.
Both drivers instead track a continuous, \emph{unwrapped} count
\begin{equation}
u_k = u_{k-1} + \mathrm{wrap}_{\pm N/2}\!\left(r_k - r_{k-1}\right),
\label{eq:unwrap}
\end{equation}
where $r_k$ is the raw reading and $N$ the counts per turn, and the wind driver averages $u$ from~\eqref{eq:unwrap} over 20 samples with a
running sum. The encoder driver originally did the same over ten samples. On
the final firmware it reads at \SI{10}{\hertz} without averaging, because its
readings proved stable and the $\pm\SI{90}{\degree}$ stops keep the wing from
reaching the seam. The
navigation code uses the circular mean
$\bar{\theta} = \operatorname{atan2}(\sum\sin\theta_i, \sum\cos\theta_i)$ for
the same reason.

The fusion task wakes on each new IMU sample. It averages the last ten BNO055
quaternions, normalizes the result, and converts it to Euler angles. Heading is
corrected from magnetic to true north with the local declination
(\SI{11}{\degree} in Salt Lake City) and for the sensor's mounting offset. The
task also estimates surge speed from successive GPS fixes.

\subsection{Sail and Heading Control}
Fig.~\ref{fig:control} shows the two control chains. Each consists of a
target source, a PID controller, an arbiter, and a servo. The PID
implementation~\cite{astrom1995} computes
\begin{align}
u_k &= \mathrm{sat}\!\left(K_p e_k + I_k - K_d\,\frac{y_k - y_{k-1}}{T}\right), \label{eq:pid}\\
I_k &= I_{k-1} + K_i T e_k , \label{eq:integral}
\end{align}
where $e_k$ is the error wrapped to $\pm\SI{180}{\degree}$, $y_k$ the measured
angle, and $T$ the sample period. In~\eqref{eq:pid}, the derivative acts on the measurement
rather than the error, so a new target does not cause a derivative kick. The
integral update~\eqref{eq:integral} is skipped whenever the output is
saturated in the direction of the error, which prevents windup, and the controller resets after a new target or a
gap in data.

\begin{figure*}[!t]
\centering
\begin{tikzpicture}[
  font=\scriptsize, >={Stealth[length=2mm]},
  b/.style={draw, rounded corners=2pt, fill=blue!6, align=center, minimum height=9mm, minimum width=19mm},
  p/.style={draw, rounded corners=2pt, fill=gray!12, align=center, minimum height=9mm, minimum width=19mm},
  sum/.style={draw, circle, inner sep=1.2pt}
]
% sail chain
\node[align=center] (vane) at (0,0) {wind vane\\$\gamma$};
\node[b] (lsp) at (2.3,0) {LSP (2 Hz)\\$\beta^{*} = \gamma \mp \alpha_z$};
\node[sum] (s1) at (4.3,0) {$+$};
\node[b] (pids) at (6.2,0) {sail PID (10 Hz)\\$K_p{=}2,\ K_i{=}K_d{=}0.5$};
\node[b] (arbt) at (9.5,0) {tail arbiter\\(20 Hz)};
\node[p] (tsv) at (12.1,0) {tail servo\\$\delta$};
\node[p] (wing) at (14.7,0) {self-trimming\\wing};
\draw[->] (vane) -- (lsp);
\draw[->] (lsp) -- node[above]{$\beta^{*}$} (s1);
\draw[->] (s1) -- node[above]{$e$} (pids);
\draw[->] (pids) -- node[above]{request} (arbt);
\draw[->] (arbt) -- (tsv);
\draw[->] (tsv) -- (wing);
\draw[->] (wing.east) -- ++(0.4,0) |- node[pos=0.75,below]{wing angle $\beta$ (AS5600)} ([yshift=-8mm]s1.south) -- node[pos=0.7,right]{$-$} (s1.south);
\draw[->] (4.3,1.05) node[above]{semi: ground station $\beta^{*}$} -- (s1.north);
\draw[->] (9.5,1.05) node[above]{manual: ground station $\delta$} -- (arbt.north);
% heading chain
\begin{scope}[yshift=-2.9cm]
\node[align=center] (src) at (0,0) {target\\$\psi^{*}$};
\node[sum] (s2) at (4.3,0) {$+$};
\node[b] (pidh) at (6.2,0) {heading PID (10 Hz)\\$K_p{=}1$};
\node[b] (arbr) at (9.5,0) {rudder arbiter\\(20 Hz)};
\node[p] (rsv) at (12.1,0) {rudder servo};
\node[p] (hull) at (14.7,0) {hull and keel};
\draw[->] (src) -- node[above]{ground station or LTP} (s2);
\draw[->] (s2) -- node[above]{$e_\psi$} (pidh);
\draw[->] (pidh) -- node[above]{request} (arbr);
\draw[->] (arbr) -- (rsv);
\draw[->] (rsv) -- (hull);
\draw[->] (hull.east) -- ++(0.4,0) |- node[pos=0.75,below]{heading $\psi$ (BNO055 and fusion)} ([yshift=-8mm]s2.south) -- node[pos=0.7,right]{$-$} (s2.south);
\draw[->] (9.5,1.05) node[above]{manual: ground station rudder angle} -- (arbr.north);
\end{scope}
\end{tikzpicture}
\caption{Sail (top) and heading (bottom) control chains. In AUTO the planners
set the targets; in SEMI the ground station sets the PID target; in MANUAL the
arbiter passes the ground-station servo angle straight through.}
\label{fig:control}
\end{figure*}

\paragraph*{Live sail planning}
The wind vane gives the apparent wind angle $\gamma$ relative to the bow
(positive from starboard), and the encoder gives the wing angle $\beta$ on the
same scale. The angle of attack is therefore $\alpha = \gamma - \beta$. The
live sail planner (LSP) chooses a wing-angle target that points the wing into
the apparent wind and then rotates it toward the bow by an angle of attack
$\alpha_z$ that depends on the point of sail:
\begin{equation}
\beta^{*} = \mathrm{wrap}_{\pm 180}\bigl(\gamma - \operatorname{sgn}(\gamma)\,\alpha_z(|\gamma|)\bigr),
\label{eq:lsp}
\end{equation}
with $\alpha_z$ in~\eqref{eq:lsp} taken from Table~\ref{tab:lsp}. A \SI{3}{\degree} deadband keeps
the planner from chasing vane noise, and it runs only when the tail arbiter is
in AUTO. The values in Table~\ref{tab:lsp} are the initial firmware settings. The
team will refine them through further water testing by logging boat speed
against angle of attack on each point of sail and keeping the angle that
gives the most forward drive for this wing's lift and drag characteristics.

\begin{table}[!t]
\caption{Sail Planner Angle of Attack by Point of Sail}
\label{tab:lsp}
\centering
\begin{tabular}{@{}lcc@{}}
\toprule
Point of sail & $|\gamma|$ (deg) & $\alpha_z$ (deg) \\
\midrule
In irons      & 0--30    & 0 \\
Close-hauled  & 30--50   & 30 \\
Close reach   & 50--75   & 35 \\
Beam reach    & 75--105  & 40 \\
Broad reach   & 105--130 & 45 \\
Training run  & 130--150 & 55 \\
Run           & 150--180 & 90 \\
\bottomrule
\end{tabular}
\end{table}

\paragraph*{Sail PID}
Rather than commanding the tail angle open loop, the sail PID closes the loop
on the measured wing angle $\beta$ and outputs the tail deflection $\delta$,
limited to $\pm\SI{90}{\degree}$. Section~\ref{sec:mech} showed that the
self-trimming equilibrium depends only on $\delta$, but the gain from tail
deflection to angle of attack depends on the tail and wing geometry and is not
known precisely; the simulator currently estimates it at
\SI{0.5}{\degree} of angle of attack per degree of tail. Measuring $\beta$
directly and integrating the error removes the steady offset that an
uncertain gain would leave. The gains are $K_p = 2$ degrees of tail per degree
of error, $K_i = 0.5$, and $K_d = 0.5$. The derivative term, acting on the
measured wing angle, damps the wing's swing as it weathervanes toward
equilibrium.

\paragraph*{Heading PID}
The heading PID wraps the heading error as
\begin{equation}
e_\psi = \bigl((\psi^{*} - \psi + \SI{180}{\degree}) \bmod \SI{360}{\degree}\bigr) - \SI{180}{\degree}
\label{eq:wrap}
\end{equation}
so the boat always turns the short way, and feeds $e_\psi$ from~\eqref{eq:wrap} to~\eqref{eq:pid}. It runs at \SI{10}{\hertz} and outputs a rudder angle limited to
$\pm\SI{35}{\degree}$. On the final firmware it is proportional only, with
$K_p = 1$ degree of rudder per degree of error; an earlier version used
$K_i = 0.05$. With no target from the ground station or planner, it holds the
heading the boat settles on five seconds after the first valid reading. If
orientation data become unavailable, it stops sending requests, and the
rudder arbiter centers the rudder after its timeout. A check that treated an
unchanging heading as a frozen IMU was disabled after it stopped the PID
during steady sailing.

\subsection{Actuator Arbiters}
Each servo has one arbiter task, the only code allowed to write its pulse
width. The arbiter's mode (MANUAL, AUTO, SEMI, or SAFE) selects which source
it obeys: a manual angle from the ground station or a request from the PID.
A PID request older than \SI{200}{\milli\second} (rudder) or
\SI{300}{\milli\second} (tail) is ignored, and the actuator moves to its safe
position: rudder centered, tail at \SI{0}{\degree}. With the tail at
\SI{0}{\degree} in line with the wing chord, \eqref{eq:moment} puts the
wing's equilibrium at zero angle of attack, so the wing should feather. This
has not yet been verified on the water; the boat did not enter SAFE mode
during the field test (Section~\ref{sec:field}). Every
output is clamped and slew-limited to \SI{6}{\degree} per \SI{50}{\milli\second}
cycle (\SI{120}{\degree\per\second}), which protects the servos and linkages
from step commands. Because each actuator has one writer and a timeout, a
crashed or starved controller leaves its actuator in a safe state rather than
frozen at its last command.

\subsection{Operating Modes and Commander}
The proposal required manual control, autonomous sailing, and semi-autonomous
sailing with either actuator under manual control. The commander task
implements these as eleven operating modes (Table~\ref{tab:modes}). Each
actuator is independently \emph{auto} (PID with a planner target),
\emph{semi} (PID with a target from the ground station), or \emph{manual}
(servo follows a ground-station angle directly). STANDBY and SAFE put both
actuators in their safe positions. Switching modes sets both arbiter modes and
suspends any PID that is not needed through FreeRTOS task notifications, so a
disabled loop always stops at the top of its cycle and restarts with fresh
history.

\begin{table}[!t]
\caption{Operating Modes}
\label{tab:modes}
\centering
\begin{tabular}{@{}lcc@{}}
\toprule
Mode & Rudder & Sail \\
\midrule
AUTO (boot default)       & auto   & auto \\
FULL\_SEMI                & semi   & semi \\
FULL\_MAN                 & manual & manual \\
AUTO\_SAIL\_MAN\_RUDDER   & manual & auto \\
AUTO\_RUDDER\_MAN\_SAIL   & auto   & manual \\
SEMI\_RUDDER\_MAN\_SAIL   & semi   & manual \\
SEMI\_RUDDER\_AUTO\_SAIL  & semi   & auto \\
SEMI\_SAIL\_MAN\_RUDDER   & manual & semi \\
SEMI\_SAIL\_AUTO\_RUDDER  & auto   & semi \\
STANDBY, SAFE             & safe   & safe \\
\bottomrule
\end{tabular}
\end{table}

The commander is the highest-priority task. It wakes on the the inter-processor communication (IPC) doorbell or
every \SI{100}{\milli\second}, decodes each command, and range-checks its
payload. Commands set the mode, drive a servo directly, set a PID target, or
upload a route. A command that does not apply in the current mode, such as a
manual rudder angle while the rudder is in AUTO, is rejected. Every command
produces a result code (for example OK, bad value, wrong mode, or bad CRC) that
is returned in telemetry with the command's sequence number, so the ground
station knows whether to resend. Gaps in the sequence are reported as dropped
commands. A compile-time option drops any manual or semi mode to SAFE after
\SI{5}{\second} without a command or heartbeat; it is currently disabled for
bench work and must be enabled before unattended operation on the water.

\subsection{Navigation}
\paragraph*{Route manager}
The ground station uploads a route as a start command, up to 32 waypoint
commands, and a commit. Each waypoint carries its index, latitude and
longitude as signed 32-bit integers in units of $10^{-7}$ degrees, and the
half-width of the corridor around the leg that ends at it. The commit carries
the waypoint count, a CRC-16/CCITT checksum over the waypoints, an arrival
radius (default \SI{8}{\meter}), and a warning margin beyond each corridor.
Waypoints accumulate in a spare buffer. The commit swaps it into the active
route, under a mutex, only if every index arrived and the checksum matches. A
failed upload therefore never disturbs the route being sailed, and a version
counter tells the planner when a new route is active.

\paragraph*{Geometry}
Over the few-kilometer scale of a mission, the navigation code projects
latitude $\phi$ and longitude $\lambda$ onto a local tangent plane around a
reference point $(\phi_0, \lambda_0)$:
\begin{equation}
n = k(\phi - \phi_0), \qquad e = k(\lambda - \lambda_0)\cos\phi_0 ,
\label{eq:local}
\end{equation}
with $k = \SI{111320}{\meter}$ per degree. Using~\eqref{eq:local}, the bearing between two points is
$\operatorname{atan2}(\Delta e, \Delta n)$. For a leg from start to end with
direction vector $\mathbf{d}$ and boat offset $\mathbf{b}$ from the start, the
signed cross-track and along-track distances are
\begin{equation}
x_{\mathrm{xt}} = \frac{d_n b_e - d_e b_n}{|\mathbf{d}|}, \qquad
x_{\mathrm{at}} = \frac{\mathbf{d}\cdot\mathbf{b}}{|\mathbf{d}|},
\label{eq:xtrack}
\end{equation}
with $x_{\mathrm{xt}} > 0$ to starboard of the path. The planner compares $x_{\mathrm{xt}}$ from~\eqref{eq:xtrack} with each leg's corridor. The route manager and
geometry library contain no RTOS or hardware calls, so they compile and run on
a laptop under a unit-test harness (Section~\ref{sec:test}).

\paragraph*{Live tack planner}
The live tack planner (LTP) turns the active route into a heading target for
the heading PID. It is implemented on the \texttt{ltp} branch and has not yet
been run on the water. It runs at \SI{2}{\hertz} as a pure function with no
RTOS calls, so the same code runs under the unit-test harness. Each cycle it
checks that GPS, IMU, and wind data are fresh and valid. It then updates a
ten-second circular mean of the wind direction, estimated as heading plus
apparent wind angle. Wind samples are skipped during a tack because the vane
lags the turn. Optionally, the planner can recover the true wind from the GPS
velocity using $\mathbf{v}_T = \mathbf{v}_A + \mathbf{v}_B$. The state
machine in Fig.~\ref{fig:fsm} follows the short-course approach of Stelzer and
Pr\"oll~\cite{stelzer2008}:
\begin{itemize}
\item \emph{DIRECT}: if the waypoint bearing is at least
\SI{45}{\degree} from the wind, steer along the leg, angled back toward the
line by up to \SI{45}{\degree} in proportion to $x_{\mathrm{xt}}$ from
\eqref{eq:xtrack}, but never into the no-go zone.
\item \emph{UPWIND\_PORT} and \emph{UPWIND\_STBD}: if the waypoint is upwind,
sail close-hauled at \SI{45}{\degree} to the wind on the tack nearer the
current heading, so the first move never turns through the wind.
\item \emph{TACKING}: tack when $|x_{\mathrm{xt}}|$ reaches 80\% of the leg's
corridor half-width while the heading is taking the boat further out, or when
the waypoint can be reached on the other tack. A tack is complete within
\SI{20}{\degree} of the new heading. Tacks are at least \SI{15}{\second}
apart, and a \SI{5}{\degree} hysteresis prevents flip-flopping at the layline.
\item \emph{ARRIVED} and \emph{FAULT}: a waypoint is reached within the
arrival radius, or once the boat passes the end of the leg inside the
corridor. Leaving the corridor by more than the warning margin latches a fault.
Both states put the boat in SAFE through the commander.
\end{itemize}
The planner sends a new heading target to the PID only when it changes by
more than \SI{3}{\degree}. A compile-time test route (an
\SI{80}{\meter} $\times$ \SI{60}{\meter} box with one leg in each direction, so
that one leg is always upwind) loads automatically at the first GPS fix for
bench and simulator runs.

\begin{figure}[!t]
\centering
\begin{tikzpicture}[
  font=\scriptsize, >={Stealth[length=2mm]},
  st/.style={draw, rounded corners=3pt, minimum width=16mm, minimum height=6.5mm, align=center, fill=blue!6},
  en/.style={draw, rounded corners=3pt, minimum width=16mm, minimum height=6.5mm, align=center, fill=gray!12}
]
\node[en] (idle) at (0,0)      {IDLE};
\node[st] (dir)  at (3.0,0)    {DIRECT};
\node[en] (arr)  at (6.3,0)    {ARRIVED};
\node[st] (up)   at (3.0,-2.6) {UPWIND\\PORT or STBD};
\node[st] (tk)   at (6.3,-2.6) {TACKING};
\node[en] (flt)  at (0,-2.6)   {FAULT};
\draw[->] (idle) -- node[above]{route loaded} (dir);
\draw[->] (dir) -- node[above]{last waypoint} (arr);
\draw[->] ([xshift=-3mm]dir.south) -- node[left,align=right]{waypoint\\upwind} ([xshift=-3mm]up.north);
\draw[->] ([xshift=3mm]up.north) -- node[right,align=left]{layline on\\same tack} ([xshift=3mm]dir.north |- dir.south);
\draw[->] ([yshift=1.5mm]up.east) -- node[above]{corridor edge} ([yshift=1.5mm]tk.west);
\draw[->] ([yshift=-1.5mm]tk.west) -- node[below]{within \SI{20}{\degree}} ([yshift=-1.5mm]up.east);
\draw[->] (up) -- node[above]{stale inputs} (flt);
\draw[->] (dir.north east) to[out=45,in=135,looseness=5] node[above]{next waypoint} (dir.north west);
\end{tikzpicture}
\caption{Live tack planner states. A fault (stale inputs, too little wind
history, or leaving the corridor) can occur from any sailing state, and
ARRIVED or a corridor fault puts the boat in SAFE.}
\label{fig:fsm}
\end{figure}

\subsection{Telemetry}
The telemetry task packs a 73-byte record into SRAM4 once per second. The
record contains pack and cell voltages; encoder angle, magnetic field
strength, and status; apparent wind angle; GPS position, altitude, satellite
count, fix quality, and time; IMU calibration state and temperature; fused
heading, roll, and pitch; both servo angles as actually driven after clamping
and slew limiting; the mode, desired heading, active leg, and cross-track
distance; the last command's sequence number and result; and uptime. Each
sensor group carries a validity byte. On each TIM6 tick the CM4 copies the
record, retrying about \SI{1}{\milli\second} later if the CM7 is mid-write,
prepends a five-byte header (destination, source, sequence number, length,
and status flags), and transmits the 78-byte frame. Commands from the ground
station use the same header with a payload of up to 16 bytes.

\subsection{Ground Station}
The ground station radio is a SparkFun SAMD21 Pro RF board running a
transparent bridge. Bytes arriving from the laptop over USB are collected
until the line has been idle for \SI{10}{\milli\second} and then sent as one
LoRa packet. Every packet received is written back to USB byte for byte. The
bridge uses the same radio settings as the boat (SF8, \SI{125}{\kilo\hertz},
CR 4/5, \SI{17}{dBm}). All protocol logic runs on the laptop, in three
programs:
\begin{itemize}
\item \texttt{gs\_term}, a terminal program in C. It validates each telemetry
frame's addresses and length and redraws a text status screen. It sends
typed commands for modes, heading and sail-angle targets, and direct servo
angles, and a heartbeat every second. It logs every raw byte with a
timestamp.
\item \texttt{gs\_web}, a browser dashboard written with only the Python
standard library. It reads the telemetry layout from the same C header at
startup, so the firmware structure stays the single definition of the
packet.
\item \texttt{plot\_log}, which converts a session log into CSV files,
twelve plots, and an HTML summary. It decodes the raw radio bytes rather than
the terminal's decoded text, so the plots show exactly what was received. The
field-test figures in Section~\ref{sec:test} were produced with it.
\end{itemize}
A simulated CM4 and bridge on a pseudo-terminal lets the ground station
software be tested without the boat.

Both radios are half duplex, which shaped the command timing. The bridge
cannot read its serial port while it transmits (about \SI{70}{\milli\second}
for a 7-byte command), so the terminal spaces packets at least
\SI{150}{\milli\second} apart. The boat cannot receive while it sends its
\SI{246}{\milli\second} telemetry frame. The team measured that any command
sent between 87 and \SI{347}{\milli\second} before a telemetry packet arrives
is lost. The terminal therefore predicts the next arrival from the boat's
measured \SI{0.997}{\second} telemetry period and holds any command that
would fall in that window.
